Code review is limited by the evidence visible to the reviewer. When an edit
breaks code in another file, neither the diff nor the reviewer's memory
guarantees that the consequence will be found. This thesis asks how repository
context changes the coverage of LLM-generated reviews, whether structural
context enables cross-file reasoning, and whether human raters perceive the
resulting difference.

A pipeline generates reviews under four modes: a diff-only baseline, a
structural arm that supplies edges from a Joern code property graph, a
retrieval arm that supplies embedding-similar code, and a hybrid containing
both. The modes are evaluated on 35 pull requests from five repositories, on
40 controlled defect injections, and in a human study.

On the pull requests, a panel of five language-model judges scores each review
on 25 criteria. The graph arm gains $+0.69$ on the nine criteria designated
in advance as structurally relevant, and 96\,\% of its total advantage falls
within them (both $p = 0.003$). Retrieval produces the largest
full-rubric gain ($+1.03$, $p = 0.002$). After correction for multiple
comparisons, the graph retains its targeted gain and retrieval its full-rubric
gain; the hybrid retains neither.

Among 28 injected defects whose consequences lie outside the edited file, the
baseline detects none, retrieval detects 5, and the graph detects 18. Adding
inheritance edges raises detection to 26 of 28. On 12 control defects
visible in the edited file, every mode
performs at or near ceiling. The eight additional detections from inheritance
edges show that graph construction has substantial room to improve.

In the human study, 27 raters compare graph and baseline reviews.
They prefer the graph review for naming affected code (mean preference
$0.710$, where $0.5$ denotes no preference; $p = 0.0005$), show no
preference on overall usefulness, and prefer the baseline review's
commentary on code clarity.

The contribution is a bounded capability, not a general improvement in review
quality. Structural context exposes cross-file consequences that a diff does
not contain; similarity retrieval improves broader rubric coverage but does
not replace structural evidence. The benefit is large when cross-file
reasoning is required and smaller but targeted on the pull-request
benchmark.
